\chapter{两种积分的比较}
\label{chap:riemann-lebesgue-comparison}

两种积分对连续函数给出同一个数，但构造方法不同。Riemann 积分要求对定义域
作有限分割，并控制每个小区间内的振荡；Lebesgue 积分允许先忽略零测集，
再用可测简单函数逼近。比较二者不只是核对一个公式：它说明 Riemann
可积性的限制恰好落在不连续点集的测度上。

\section{上下阶梯函数}

设 \(f:[a,b]\to\mathbb R\) 有界，分割
\(P:a=x_0<\cdots<x_n=b\)。记
\[
 m_k=\inf_{[x_{k-1},x_k]}f,\qquad
 M_k=\sup_{[x_{k-1},x_k]}f.
\]
忽略有限个分点处的取值，定义下、上阶梯函数
\[
 s_P=\sum_{k=1}^{n}m_k\mathbf1_{(x_{k-1},x_k)},\qquad
 t_P=\sum_{k=1}^{n}M_k\mathbf1_{(x_{k-1},x_k)}.
\]
它们是 Lebesgue 可测简单函数，并且几乎处处满足
\(s_P\le f\le t_P\)。此外
\[
 \int s_P\,dm=L(f,P),\qquad
 \int t_P\,dm=U(f,P),
\]
其中右侧是 Darboux 下和、上和。有限个分点的选择不影响 Lebesgue 积分。

\begin{theorem}[两种积分的一致性]
\label{thm:riemann-lebesgue-agreement}
若有界函数 \(f:[a,b]\to\mathbb R\) Riemann 可积，则 \(f\) Lebesgue
可测且可积，并且
\[
 \int_a^b f(x)\,dx_{\mathrm R}
 =\int_{[a,b]}f\,dm .
\]
\end{theorem}

\begin{proof}
由 Riemann 可积性，可取分割 \(P_n\) 使
\[
 U(f,P_n)-L(f,P_n)<2^{-n}.
\]
用前 \(n\) 个分割的公共加细代替 \(P_n\)，可设分割递细。令 \(D\) 为
这些分割点的并；它至多可数。在 \([a,b]\setminus D\) 上，相应阶梯函数
满足 \(s_n\uparrow\)、\(t_n\downarrow\)，且
\[
 \int(t_n-s_n)\,dm\longrightarrow0.
\]
在 \([a,b]\setminus D\) 上令 \(s=\sup_ns_n\)、\(t=\inf_nt_n\)，并在
\(D\) 上把二者都定义为零。它们是可测函数，并且
\(s\le f\le t\) 在 \([a,b]\setminus D\) 上成立。由 Fatou 引理，
\[
 0\le\int(t-s)\,dm
 \le\liminf_n\int(t_n-s_n)\,dm=0.
\]
所以 \(f=s=t\) 几乎处处，因而 \(f\) 可测。又 \(|f|\) 有界且区间
测度有限，故 \(f\in L^1\)。若 \(|f|\le M\)，则
\(|s_n|,|t_n|\le M\) 几乎处处；控制收敛定理给出
\[
 L(f,P_n)=\int s_n\longrightarrow\int f,\qquad
 U(f,P_n)=\int t_n\longrightarrow\int f,
\]
两端同时趋于 Riemann 积分值。
\end{proof}

\section{局部振幅与不连续点}

对有界函数 \(f:[a,b]\to\mathbb R\)，定义点 \(x\) 处的振幅
\[
 \omega_f(x)
 =\inf_{\delta>0}
 \left(
 \sup_{y\in[a,b]\cap(x-\delta,x+\delta)}f(y)
 -\inf_{y\in[a,b]\cap(x-\delta,x+\delta)}f(y)
 \right).
\]

\begin{proposition}
\label{prop:oscillation-continuity}
\(f\) 在 \(x\) 连续当且仅当 \(\omega_f(x)=0\)。对每个
\(\varepsilon>0\)，集合
\[
 D_\varepsilon=\{x\in[a,b]:\omega_f(x)\ge\varepsilon\}
\]
是相对闭集。因此不连续点集
\[
 D(f)=\bigcup_{r=1}^{\infty}D_{1/r}
\]
是 \(F_\sigma\) 集。
\end{proposition}

\begin{proof}
连续性直接给出任意小的邻域振幅。反之，若邻域内函数值的振幅小于
\(\varepsilon\)，则其中任意 \(y\) 满足
\(|f(y)-f(x)|<\varepsilon\)。

若 \(x\notin D_\varepsilon\)，存在邻域 \(U\) 使 \(f\) 在 \(U\) 上的
振幅小于 \(\varepsilon\)。对充分靠近 \(x\) 的 \(z\)，可选取包含于
\(U\) 的 \(z\) 的邻域，故 \(\omega_f(z)<\varepsilon\)。所以
\([a,b]\setminus D_\varepsilon\) 相对开。
\end{proof}

\section{Lebesgue 判别准则}

\begin{theorem}[Lebesgue 的 Riemann 可积性判别]
\label{thm:lebesgue-riemann-criterion}
有界函数 \(f:[a,b]\to\mathbb R\) Riemann 可积，当且仅当其不连续点集
\(D(f)\) 的 Lebesgue 测度为零。
\end{theorem}

\begin{proof}
先设 \(f\) Riemann 可积。固定 \(\varepsilon>0\)，取分割 \(P\) 使
\[
 \sum_{I\in P}\omega(f;I)|I|<\eta,
\]
其中 \(\omega(f;I)=\sup_I f-\inf_I f\)。除有限个分点外，若
\(x\in D_\varepsilon\)，则 \(x\) 所在分割小区间的振幅至少为
\(\varepsilon\)。这些坏区间的总长度至多 \(\eta/\varepsilon\)，故
\(m^*(D_\varepsilon)\le\eta/\varepsilon\)。令 \(\eta\downarrow0\)，
得到 \(m(D_\varepsilon)=0\)。由
\(D(f)=\bigcup_{r=1}^{\infty}D_{1/r}\) 得 \(m(D(f))=0\)。

反过来，设 \(m(D(f))=0\)，并令 \(|f|\le M\)。若 \(M=0\)，函数恒为零，
结论成立。以下设 \(M>0\)。给定
\(\varepsilon>0\)，选 \(r\) 充分大，使
\((b-a)/r<\varepsilon/2\)。闭集 \(D_{1/r}\) 零测，可由有限个开区间
覆盖，且总长度小于 \(\varepsilon/(4M)\)。记这些区间的并为 \(G\)，并令
\[
 K=[a,b]\setminus G.
\]
对每个 \(x\in K\)，有 \(\omega_f(x)<1/r\)，故存在相对开邻域 \(U_x\)，
使 \(f\) 在 \(U_x\) 上的振幅小于 \(1/r\)。从 \((U_x)_{x\in K}\) 中取
有限子覆盖。它们与 \(G\) 一起构成 \([a,b]\) 的有限开覆盖。由 Lebesgue
数引理，可取分割 \(P\)，使每个分割小区间的闭包包含在 \(G\) 或某个
\(U_x\) 中。前一类小区间的总长度不超过 \(m(G)\)，其振幅至多 \(2M\)；
后一类小区间的振幅小于 \(1/r\)。于是
\[
 U(f,P)-L(f,P)
 \le 2M m(G)+\frac1r(b-a)
 \le 2M\sum_j|I_j|+\frac1r(b-a)<\varepsilon.
\]
由 Darboux 判据，\(f\) Riemann 可积。
\end{proof}

证明中“有界”不能省略。函数 \(x\mapsto x^{-1/2}\) 在 \((0,1]\) 连续，
不连续点（若补定义 \(f(0)\)）只有一个，却不在通常有界 Riemann 积分的
定义域内；它可作为反常 Riemann 积分和 Lebesgue 积分处理。

\section{典型函数}

\begin{example}[Dirichlet 函数]
\[
 f=\mathbf1_{\mathbb Q\cap[0,1]}
\]
处处不连续，故不 Riemann 可积；它是可测函数并且 Lebesgue 积分为零。
\end{example}

\begin{example}[Thomae 函数]
令
\[
 t(x)=
 \begin{cases}
 1/q,&x=p/q\in[0,1],\ (p,q)=1,\ q\ge1,\\
 0,&x\notin\mathbb Q.
 \end{cases}
\]
它在每个无理点连续，在每个有理点不连续。不连续点集可数，故它 Riemann
可积，积分为零。这里不能只说“有理数零测”：还必须先验证无理点连续。
\end{example}

\begin{proof}
固定无理点 \(x\) 和 \(\varepsilon>0\)。使 \(1/q\ge\varepsilon\) 的
约分有理数在 \([0,1]\) 中只有有限个。取 \(x\) 的小邻域避开这些点，
邻域内便有 \(t(y)<\varepsilon\)。在有理点 \(p/q\)，任意邻域含无理点，
函数值可趋于零，与 \(t(p/q)=1/q\) 不相容。
\end{proof}

\begin{example}[稠密开集的指示函数]
若 \(G\subset[0,1]\) 是稠密开集，且其补集也无孤立点，则
\(\mathbf1_G\) 在 \(\partial G=[0,1]\setminus G\) 上不连续，在 \(G\)
内连续。因此它 Riemann 可积当且仅当 \(m(\partial G)=0\)，积分等于
\(m(G)\)。取 \(G\) 为 Cantor 集的补集，可得到不连续点集不可数的例子；
前面的 Thomae 函数则给出不连续点集稠密而零测的例子。
\end{example}

\section{有界与无界积分的分界}

在有限区间上，任何有界可测函数都 Lebesgue 可积，而 Riemann 可积还要求
不连续点集零测。无界函数则需另外检查 \(\int|f|<\infty\)。例如
\[
 f(x)=x^{-\alpha}\quad(0<x\le1)
\]
在 Lebesgue 意义下可积当且仅当 \(\alpha<1\)，这与反常 Riemann 积分的
绝对收敛条件一致。Lebesgue 积分不会自动给条件收敛的反常积分赋值：
\(\int f\) 作为有限实数存在等价于 \(\int|f|<\infty\)。

\section*{习题}
\addcontentsline{toc}{section}{习题}
\subsection*{A组　两种积分与振幅}
\begin{enumerate}[label=\textbf{\thechapter.\arabic*.}]
\exerciseitem{riemann-step-integrals}
验证上下阶梯函数的 Lebesgue 积分分别等于 Darboux 下和、上和，并说明
分点处怎样取值均不影响结论。

\exerciseitem{riemann-measurability-direct}
设 \(f\) Riemann 可积。利用上下阶梯函数列直接证明 \(f\) 是 Borel
可测函数经过零测集修改所得的 Lebesgue 可测函数。

\exerciseitem{oscillation-upper-semicontinuity}
证明 \(x\mapsto\omega_f(x)\) 是上半连续的，即
\(\{\omega_f<\varepsilon\}\) 为开集。

\exerciseitem{riemann-criterion-forward-detail}
补全定理 \ref{thm:lebesgue-riemann-criterion} 正向证明中分点和区间端点
的处理。

\end{enumerate}
\subsection*{B组　判别准则与典型函数}
\begin{enumerate}[label=\textbf{\thechapter.\arabic*.},resume]
\exerciseitem{riemann-criterion-backward-detail}
补全反向证明中的有限分割构造，并给出只依赖 \(M,b-a,\varepsilon\) 的
上和减下和估计。

\exerciseitem{thomae-continuity-set}
完整求出 Thomae 函数的连续点集和不连续点集，并分别用 Riemann 理论与
Lebesgue 理论计算其积分。

\exerciseitem{indicator-riemann-boundary}
对 \(E\subset[a,b]\)，证明 \(\mathbf1_E\) 的不连续点集恰为
\(\partial E\)。由此证明 \(\mathbf1_E\) Riemann 可积当且仅当
\(m(\partial E)=0\)。

\exerciseitem{fat-cantor-indicator}
设 \(F\subset[0,1]\) 是正测度、无内点闭集。判断 \(\mathbf1_F\) 的
Riemann 可积性与 Lebesgue 积分，并解释二者差异。

\end{enumerate}
\subsection*{C组　反例与构造}
\begin{enumerate}[label=\textbf{\thechapter.\arabic*.},resume]
\exerciseitem{dense-zero-discontinuities}
构造一个 Riemann 可积函数，其不连续点集在 \([0,1]\) 中稠密。

\exerciseitem{improper-lebesgue-comparison}
对 \(f(x)=x^{-\alpha}\sin(1/x)\)，讨论它在 \((0,1]\) 上的反常 Riemann
积分和 Lebesgue 可积性，区分 \(\alpha<1\)、\(1\le\alpha<2\) 与
\(\alpha\ge2\)。

\optionalexerciseitem{bounded-ae-continuous-version}
设 \(f\) 是 \([a,b]\) 上的有界 Lebesgue 可测函数，且几乎处处连续。
证明 \(f\) Riemann 可积。说明“几乎处处等于某个 Riemann 可积函数”并不
蕴含 \(f\) 自身 Riemann 可积。

\projectexerciseitem{prescribed-discontinuity-set}
证明：对每个零测 \(F_\sigma\) 集 \(E\subset[a,b]\)，存在有界 Riemann
可积函数，其不连续点集恰为 \(E\)。
\end{enumerate}
